% P4 executed GitHub Action 38315929126; synthetic source only.
\subsection{Does simply adding more shrinkage intervals close the gap?}
\label{sec:p4-refinement}
To test whether the observed benefit of the
whole-numerator Bernstein formulation is merely
an artefact of using insufficiently many intervals,
we reran the same fifteen deterministic synthetic
source spectra with $12$, $24$ and $48$
equal-length subintervals. The original
$12$-subinterval run already used source
pair features held fixed across comparisons.
The measured mean improvement of Bernstein
$48$ segments over its own $12$-segment
coverage was just \emph{$+0.0412$ percentage
points} over these fifteen cases, with a
strict increase in only $2/15$.
No case gave a decrease beyond the
stated numerical tolerance in this test.

The result supports a narrow, mechanistic
reading of the earlier synthetic comparisons:
the much larger differences between V23
termwise/Taylor and V25 whole-numerator
Bernstein at identical $12$-segment budgets
cannot be attributed exclusively to
insufficient segmentation of the latter
bound. Exploiting polynomial cancellations
is a different enclosure choice from
simply refining the grid. This does not
prove optimality of Bernstein bounds,
because neither enclosure includes exact
root isolation and each is conditional
on floating-point eigensystems.
The $12$-, $24$- and $48$-segment runs also
do not establish better runtime efficiency;
denser subdivision has a computational cost.
All cases are \emph{synthetic}, not
additional real EEG recordings.
